\subsection{Rational adapted normal forms and the original labels}
\label{sec:flag-adapted-normal-form}

Let \(E\) be the original CM host, with conjugation denoted
by a bar. Use a nondegenerate hermitian form \(B\) linear
in its first variable. Write the two rational flag bases as
\(e_0,\ldots,e_{r-1},z=e_r\) and
\(f_0,\ldots,f_{r-1},z'=f_r\).
Let \(X_i,Y_i\) be the spans of the first \(i\) isotropic
vectors, \(X=X_r,Y=Y_r\), and \(I=X\cap Y\).
The last vectors are perpendicular to their respective
isotropic spaces and have the same norm \(-T\), \(T>0\).
The original positivity hypothesis says, at each complex
embedding, that a vector perpendicular to \(X,z\) and not
in \(X\) has positive norm, and similarly for \(Y,z'\).
In particular an isotropic vector perpendicular to \(X,z\)
belongs to \(X\).

An adapted change replaces each vector by itself plus a
combination of the preceding isotropic vectors (including
the change of \(z\) by a vector in \(X\)).
For a skew-hermitian endomorphism \(A\), the flag Lie
condition is \(Ae_i\in X_i\) for \(i<r\) and
\(A(X^\perp)\subset X\). Its simple coefficients are the
coefficients of \(e_i\) in \(Ae_{i+1}\), modulo \(X_i\),
for \(0\le i<r\). Relevance means that the traces to the
totally real subfield of the sums of these coefficients
for the two flags agree whenever both Lie conditions hold.
These conditions and coefficients are unchanged by adapted
changes: the flags are unchanged, and conjugation by an
upper unitriangular matrix on the span of the flag vectors
preserves the first superdiagonal of a strictly upper
triangular matrix.

Write \(D_l=\det(B(e_a,f_b))_{0\le a,b<l}\), with \(D_0=1\).
The zero-Gram condition in the source is
\[
 D_{r+1}\overline{D_{r+1}}-T^2D_r\overline{D_r}=0.
\]
Both this condition and every leading minor are invariant
under adapted changes, since the leading matrices are
multiplied by unitriangular matrices.

\begin{lemma}[The common tail]
\label{lem:rational-flag-common-tail}
Under the stated positivity, relevance and zero-Gram
conditions, put \(j=r-\dim_E I\). Then \(D_j\ne0\);
there are adapted bases with
\(f_b=c e_b\) for every \(j\le b\le r\), where
\(c\overline c=1\). The joint span has radical \(I\).
\end{lemma}

\begin{proof}
Neither last vector belongs to \(X+Y\). Indeed, writing
\(z=x+y\) would give \(B(y,x)=0\) since \(z\perp X\);
then \(B(z,z)=0\), a contradiction. Interchange the
flags for the other assertion.

Suppose first that \(I\ne0\). For a nonzero \(u\in I\)
write
\[
 u\equiv a e_i\pmod{X_i},\qquad
 u\equiv a'f_{i'}\pmod{Y_{i'}},\qquad aa'\ne0.
\]
For any \(v\perp X_{i+1}+Y_{i'+1}\), the endomorphism
\[
 A_v(w)=B(w,v)u-B(w,u)v
\]
is skew-hermitian and satisfies both flag Lie conditions.
On the first flag its only possibly nonzero simple
coefficient is \(a B(e_{i+1},v)\); on the second it is
\(a'B(f_{i'+1},v)\). Apply relevance also with every
scalar multiple of \(v\). Nondegeneracy of the trace
pairing of \(E\) over its real subfield and then
nondegeneracy of \(B\) give
\[
 a e_{i+1}-a'f_{i'+1}\in X_{i+1}+Y_{i'+1}.
\]
If both successors are isotropic, separating the two
prefix terms gives a common vector with successor levels
and the same leading coefficients \(a,a'\). If just one
successor is the last vector, that vector belongs to
\(X+Y\), which is impossible. Iterating proves
\(i=i'\) and \(z'-c z\in X+Y\) with \(c=a/a'\).
This ratio is unique because \(z\notin X+Y\).

Let \(j_0\) be the minimum leading index of a nonzero
vector of \(I\). Successor propagation supplies common
vectors at every level \(j_0,\ldots,r-1\), and no vector
of \(I\) has lower level in either flag. The successive
one-dimensional graded quotients show that
\(\dim I=r-j_0\), so \(j_0=j\), and
\(X=X_j\oplus I,\ Y=Y_j\oplus I\).
Choose the common vectors with first leading coefficient
one and second leading coefficient \(1/c\), and replace
the tail by them and their \(c\)-multiples, respectively.
Write \(z'-c z=x+y\) with \(x\in X,y\in Y\). Replace
\(z,z'\) by \(z+c^{-1}x,z'-y\); they now differ by the
factor \(c\). Their norms did not change, hence
\(c\overline c=1\).
A vector of \(X_j\) perpendicular to \(Y_j\) is
perpendicular to all of \(Y\) and the adapted \(z'\).
By positivity it lies in \(I\), hence is zero.
Thus \(D_j\ne0\).

If \(I=0\), put \(j=r\). If \(D_r=0\), the zero-Gram
identity forces \(D_{r+1}=0\). There is then a nonzero
vector \(y'\in Y+Ez'\) perpendicular to \(X,z\).
It cannot lie in \(Y\), since it would be isotropic
and belong to \(I\). On the other hand \(D_r=0\)
supplies a nonzero \(x\in X\) perpendicular to \(Y\).
One has \(B(x,z')\ne0\) by positivity for the second
flag and \(I=0\). The nonzero coefficient of \(z'\)
in \(y'\) contradicts \(B(x,y')=0\).
Consequently \(D_r\ne0\) (also for \(r=0\)).
The space \(X+Y\) is nondegenerate; adapt the last
vectors by subtracting their projections onto it.
These projections lie respectively in \(X,Y\).
The new last vectors lie in its orthogonal complement.
Expanding the last row and column of the cross
determinant gives
\(B(z,z')\overline{B(z,z')}=T^2\).
Write \(z'=cz+w\), \(w\perp z\), in that complement.
The equality forces \(c\overline c=1\) and \(B(w,w)=0\).
Since \(w\perp X,z\) and \(I=0\) with
\(X\cap(X+Y)^\perp=0\), positivity gives \(w=0\).

In both cases the adapted joint span is the orthogonal
sum of \(X_j+Y_j\), \(Ez\) and \(I\); the first two
summands are nondegenerate and the third is radical.
Its dimension is \(r+j+1\). For any vector perpendicular
to \(X,z\), its projection onto \(X_j+Y_j\) lies in
\(X_j\); subtracting it gives a vector perpendicular
to the joint span. The symmetric assertion holds for
\(Y,z'\). These also give the remaining joint-span
orthogonality properties of the rational normal form.
\end{proof}

\begin{lemma}[The signed block matrix]
\label{lem:rational-flag-signed-block}
The two leading bases of length \(j\) can be further
adapted so that their cross matrix is a block diagonal
matrix of nonzero constant anti-diagonals, multiplied
in column \(b\) by \((-1)^{r-b}\).
\end{lemma}
\begin{proof}
Use the common-tail bases just constructed and let
\(Z=(\overline{B(e_a,f_b)})_{a,b<j}\); this conjugation
is the source cross-matrix convention. For strictly
upper triangular \(j\)-by-\(j\) matrices \(A,A'\)
satisfying
\[
 A^*Z+ZA'=0,\qquad A^*=\overline A^{\,t},
\]
act by \(A,A'\) on \(X_j,Y_j\) in their respective
bases, and by zero on \((X_j+Y_j)^\perp\).
The identity is exactly the skew-hermitian condition
on cross terms. This operator satisfies both flag
Lie conditions: for instance the projection of
\(X^\perp\) onto \(X_j+Y_j\) lies in \(X_j\).
The other simple coefficients, including those at
the common tail, are zero. If \(s(A)\) denotes the
sum of the first superdiagonal entries, relevance
therefore gives equality of the traces of
\(s(A)\) and \(s(A')\).
The same holds for \(dA,\overline d A'\) for every
\(d\in E\), so the trace pairing gives
\(s(A)=\overline{s(A')}\).

By elimination write \(Z=LPU\), with \(L,U\) lower
and upper unitriangular, and \(P\) monomial:
\(P_{\sigma(b),b}=t_b\ne0\) for a permutation
\(\sigma\). Explicitly, in the first column choose
the first nonzero entry, eliminate below it by lower
row operations and to its right by upper column
operations, then repeat on the remaining ordered
rows and columns. Invertibility ensures a pivot at
every step; the earlier cleared pivot rows and columns
are preserved. The identity just proved passes to \(P\):
conjugate \(A\) by \(L^*\) and \(A'\) by \(U\) to
transport the equation, noting that unitriangular
conjugation preserves \(s\).

If \(c_0<b\) and \(\sigma(c_0)=\sigma(b)+1\), use
\(A=E_{\sigma(b),\sigma(c_0)}\) and
\(A'=-(t_b/t_{c_0})E_{c_0,b}\), where \(E_{a,b}\)
is the matrix unit. They satisfy \(A^*P+PA'=0\).
The equality \(1=\overline{s(A')}\) forces
\(b=c_0+1\) and \(t_b=-t_{c_0}\).
Thus every descent of the inverse permutation
\(\sigma^{-1}\) is by one. Its maximal descending
runs are disjoint consecutive intervals; successive
runs must have increasing intervals because their
boundary is an ascent. Hence the permutation
reverses consecutive blocks, with no other
permutation of the blocks. Within each block the
\(t_b\) alternate signs, so
\((-1)^{r-b}t_b\) is constant there.
Finally change the leading bases by the upper
unitriangular matrices
\((\overline L^{-1})^t,U^{-1}\), respectively.
The cross matrix in the convention of \(B\) is now
\(\overline P\), with the stated column signs.
The common vectors and last vectors are untouched.
The empty leading matrix is allowed.
\end{proof}

\begin{theorem}[Actual adapted normal-form incidence and labels]
\label{thm:flag-original-nf-labels}
For every term of the original relevant zero-Gram
set the rational adapted normal form exists with
the original \(j,T,\mathfrak d\), and every adapted
normal form has the original label, \(\Delta\) and
\(D_b\). These are the original assertions
\texttt{NFExC} and \texttt{LabC}.
\end{theorem}
\begin{proof}
The actual relevant-set and representative
conditions supply relevance of the rational flags,
the zero-Gram equality, \(T=T_u>0\), and the positivity
at every embedding. The rational flags are the
ones attached to the original term, not representatives
of a different double coset. The preceding two
lemmas give adapted bases, the common scalar of
norm one and the required signed block matrix.
The host conditions supply the original nonzero,
conjugation-stable ideal \(\mathfrak d\).

To complete the normal form, in the nondegenerate
orthogonal complement of \(X_j+Y_j+Ez\), choose
vectors \(d_b\), \(j\le b<r\), satisfying
\(B(e_a,d_b)=\delta_{ab}\) for the common tail.
Existence follows by nondegeneracy and independence.
Replace each \(d_b\) by
\[
 d_b-\tfrac12\sum_{a=j}^{r-1} B(d_b,d_a)e_a.
\]
Their mutual pairings become zero, while the prescribed
pairings with the tail are unchanged. They are
perpendicular to both leading bases and the last
vector. This is exactly the auxiliary rational dual
family of the original normal form (take zero in
the unused indices). No integral splitting or extension
of the host field is required in this construction.

For an arbitrary adapted normal form the isotropic
spaces and leading minors are unchanged. Its
nondegenerate leading matrix and common tail give
\(j=r-\dim I\), the original \(j_T\).
The norm of the last vector gives \(T=T_u\).
For \(0\le l\le j\), the original label is
\(\overline{D_l}\), because the original cross
matrix is conjugate to the one used for \(B\).
The normal-form normalized minor is
\[
 n_l=|N_{E/\mathbb Q}(D_l)|\,
       N(\mathfrak d)^{-2\sum_{a=0}^{l-1}(r-a)} .
\]
The original \(\operatorname{frameExp}(r,l)\) is the
displayed sum, and \(\operatorname{Nloc}\) is the
absolute norm of the principal fractional ideal
times \(\mathfrak d^{-2\operatorname{frameExp}(r,l)}\).
Conjugation preserves the field norm. Thus \(n_l\)
is exactly \(\operatorname{Nloc}\) of the original
label, and is nonzero exactly when that label is
nonzero. In particular \(n_j=\Delta\); the product
over the nonzero labels for \(1\le l\le j\) is
exactly \(D_b\), with no factor at \(l=0\).
Both truncated label functions are zero above \(j\);
below or at \(j\) they are conjugates of each other.
This proves every required label equality, including
when \(j=0\).
\end{proof}

The rational incidence and its labels are now
available in the weighted fixed-row comparison.
They do not themselves imply the original congruence
between the two row representatives.
